\section{Introduction}
\label{sec:introduction}

Early differentiation of uveal melanocytic lesions is clinically decisive: uveal
melanoma is malignant and time-sensitive, yet on fundus imaging it is hard to
separate from benign nevus or congenital hypertrophy of the retinal pigment
epithelium (CHRPE), and expert reading is subject to high inter-observer
variability \cite{chadwick2025ai}. Deep learning provides reproducible decision support, and recent
work attains strong four-class accuracy on ultra-widefield fundus images
\cite{tasso2026lesion}. These systems, however, are built on large convolutional
backbones and assume workstation- or cloud-class compute, which sits awkwardly
with the point-of-care setting in which screening is most valuable.

Bringing such a classifier onto an embedded FPGA would place inference at the
bedside, at low power and latency and without transferring sensitive images
off-device. The obstacle is one of scale. Accurate models in this domain rely on
ResNet-class networks whose weight and activation memory exceeds the on-chip
budget of a small device, where block RAM is the binding resource. Compression and low-bit quantization are therefore
mandatory, yet accuracy cannot be traded away freely: a missed melanoma carries a
direct clinical cost, so a deployment is useful only if it preserves diagnostic
quality. The problem is thus to locate the operating point, in model size and
arithmetic precision, that fits the device while keeping accuracy clinically
acceptable, a balance that the strong class imbalance of the data makes harder
still.

We address this tradeoff with three contributions:
\begin{itemize}
\item a teacher-assistant knowledge-distillation chain of decreasing capacity that
  compresses the network, where an intermediate assistant suppresses the
  degradation incurred when a large teacher is distilled directly into a small
  student \cite{hinton2015distilling, cho2019efficacy, mirzadeh2020takd};
\item a Brevitas-based study of quantization-aware training across multiple
  weight and activation bit-widths \cite{pappalardo2023brevitas};
\item an end-to-end deployment through the FINN dataflow toolchain
  \cite{umuroglu2017finn, blott2018finnr} in a resource-driven design-space
  exploration over arithmetic precision and layer width, the axes that raise
  accuracy as they grow and free on-chip memory as they shrink.
\end{itemize}

On the Ultra96, the design-space exploration produces two deployable 6-bit
accelerator families rather than a single winner. One operating point preserves
more model capacity and favors accuracy, but remains memory-bound; the other
trades a small accuracy margin for enough on-chip memory headroom to add
parallelism and improve throughput. Both deployed accelerators remain bit-exact
with their software references; the fastest point reaches 12.71 images/s at
187.5\,MHz on the 133-image test set.
